% >>> typage-template ieee-conference@2026.10.1 variant=letter
\documentclass[conference]{IEEEtran}
% IEEEtran selects Times (ptm); under XeTeX that needs the T1 encoding.
\usepackage[T1]{fontenc}
\bibliographystyle{IEEEtran}
% <<< typage-template
\usepackage{cite}
\usepackage{amsmath,amssymb,amsfonts}
\usepackage{graphicx}
\usepackage{booktabs}
\usepackage{textcomp}
\usepackage{xcolor}

\begin{document}

\title{Paper Title}

\author{\IEEEauthorblockN{First Author}
\IEEEauthorblockA{\textit{Department} \\
\textit{University}\\
City, Country \\
first.author@example.com}
\and
\IEEEauthorblockN{Second Author}
\IEEEauthorblockA{\textit{Laboratory} \\
\textit{Institute}\\
City, Country \\
second.author@example.com}}

\maketitle

\begin{abstract}
State the problem, the approach and the main result. Do not use citations or
symbols in the abstract.
\end{abstract}

\begin{IEEEkeywords}
first keyword, second keyword, third keyword
\end{IEEEkeywords}

\section{Introduction}
Introduce the problem and list the contributions. Cite prior work by
number~\cite{lamport1994latex,knuth1984texbook}.

\section{Method}
Describe the method. Number the equations you refer to:
\begin{equation}
  \label{eq:objective}
  \min_{\theta} \; \frac{1}{n} \sum_{i=1}^{n} \ell\bigl(f_\theta(x_i), y_i\bigr) .
\end{equation}

\section{Experiments}
Table~\ref{tab:results} summarizes the results of~\eqref{eq:objective}.

\begin{table}[htbp]
  \caption{Results on the Example Benchmark}
  \label{tab:results}
  \centering
  \begin{tabular}{lcc}
    \toprule
    Method   & Accuracy (\%) & Time (s) \\
    \midrule
    Baseline & 81.2          & 12.4     \\
    Ours     & 84.7          & 10.9     \\
    \bottomrule
  \end{tabular}
\end{table}

\section{Conclusion}
Conclude and outline future work.

\section*{Acknowledgment}
Acknowledge funding and help here.

\bibliography{refs}

\end{document}
